\section{Опоздание в магнитном кольце}
\label{sec:mag}

Магнитная сила перпендикулярна скорости, поэтому между зазорами модуль скорости постоянен вдоль всей траектории. Период $T=\oint\mathrm{d}l/v=C/v$, и
\begin{equation}
\label{eq:mag-Tratio}
\frac{T}{T_s}=\frac{C}{C_s}\,\frac{v_s}{v}
\end{equation}
точно. Скорость задаётся импульсом: из $\beta\gamma=\beta_s\gamma_s(1+\delta)$ и $\gamma^2=1+(\beta\gamma)^2$
\begin{equation}
\label{eq:mag-v}
\frac{v_s}{v}=\frac{\gamma(\delta)}{\gamma_s(1+\delta)}
=1-\frac{\delta}{\gamma_s^2}+\Biggl(\frac{1}{\gamma_s^4}+\frac{3\beta_s^2}{2\gamma_s^2}\Biggr)\delta^2+O(\delta^3).
\end{equation}
Слагаемое $1/\gamma_s^4$ --- квадрат линейного члена, $3\beta_s^2/(2\gamma_s^2)$ --- вторая производная скорости по импульсу; оба одного порядка. Произведение~\eqref{eq:mag-v} и~\eqref{eq:set-C} даёт коэффициенты~\eqref{eq:set-tau}:
\begin{equation}
\label{eq:mag-eta}
\eta_0=\alpha_0-\frac{1}{\gamma_s^2},
\qquad
\eta_1=\alpha_1-\frac{\alpha_0}{\gamma_s^2}+\frac{1}{\gamma_s^4}+\frac{3\beta_s^2}{2\gamma_s^2}.
\end{equation}
